\section{Purpose}

\textbf{[STIPULATION]} This document formulates the \emph{complete
translatability} between the FFGFT mode formalism and standard Hilbert-space
quantum mechanics. The translation is not an approximation and not a
specialisation, but a \textbf{bijective structural statement}: every QM
computation is mappable into FFGFT language, and every FFGFT statement about
quantum systems is reproducible in Hilbert-space language.

\textbf{Layer note.} Alongside the algebraic bridges [B] and geometric core
statements [K], this document contains a substantial proportion of conceptual
and ontological framings [S]. The [S] statements concern interpretational
questions (Born rule, determinism, Bloch-sphere position) -- they are correctly
marked, but the reader should know: the ratio here is K~:~S~$\approx$~1~:~5.
The established core is narrow; the rest is framing. That is not a deficiency
-- framings belong in a bridge document. But bridge and core are not the same.

This is the methodological point: FFGFT does not \emph{replace} quantum
mechanics but \emph{extends} it. Both formalisms deliver the same measurable
results -- up to a $\xi$-correction of order $10^{-5}$, which currently lies
below the NISQ noise floor (A160). They differ only in the ontological reading
of some basic concepts. One who computes with Hilbert-space tools can continue
to do so; the translation table below is the bridge.

The canonical Hilbert space of the full theory is
\[
  \boxed{\;\mathcal{H} = L^2(T^4)\otimes\mathbb{C}^3\;}
\]
-- square-integrable functions on the 4-torus, tensored with the three-element
$Z_3$ fibre space (A030). This document builds it in two directions: \emph{downward}
(FFGFT $\to$ standard Hilbert space, with reductions) and \emph{upward} (which
extensions the standard Hilbert space needs to natively carry the full FFGFT).

\section{Where We Are Geometrically}

\textbf{[B]} Before the translation tables, the geometric picture. The full
FFGFT geometry is the 4-torus $T^4 = T^3\times S^1$ (A030): three spatial
winding directions and one temporal. For a single qubit -- a two-level system
at fixed time, without mass dynamics -- most of these directions are inactive.
Two rotational degrees of freedom remain: one for the position between $|0\rangle$
and $|1\rangle$ (this becomes $z$), one for the phase (this becomes $\theta$).
Their natural space is the 2-torus $T^2 = S^1\times S^1$.

Pulling one major radius large, $T^2$ becomes locally the cylinder
$\mathbb{R}\times S^1$: one loop is cut, $z$ becomes a \emph{linear} axis.
Then contracting the two poles to points, the cylinder becomes the Bloch sphere
$S^2$ of standard QM. This creates a hierarchy of reductions:

\begin{center}
\renewcommand{\arraystretch}{1.3}
{\small
\begin{tabular}{lcl}
\toprule
\textbf{Geometry} & \textbf{Dim.} & \textbf{Use} \\
\midrule
4-torus $T^4$ & 4 & full FFGFT (masses, $\alpha$) \\
3-torus $T^3$ & 3 & mode space at fixed time \\
2-torus $T^2$ & 2 & single qubit, without reduction \\
Cylinder $\mathbb{R}\times S^1$ & 2 & limit $R\to\infty$, local approximation \\
Bloch sphere $S^2$ & 2 & poles contracted, standard QM \\
\bottomrule
\end{tabular}
}
\renewcommand{\arraystretch}{1.0}
\end{center}

\textbf{[B]} \textbf{What the reductions cost -- and where the loss really sits.}
Each level loses something: $T^4\to T^3$ the time winding and hence the masses;
$T^3\to T^2$ a spatial winding and hence the third generation; $T^2\to$ cylinder
\emph{one compact loop} -- the weakest reduction, and precisely here sits the
fractal correction factor $K_{\text{frak}} = 1-100\xi \approx 0{.}9867$ (A040):
$K_{\text{frak}}$ is the cylinder's memory of its torus origin. The loss is
\emph{not} the subsequent SI conversion: the duality-carrying decompactification
$S^1_m\to\mathbb{R}_t$ ($\tilde T\cdot m=1$, A060/A090) is a mode-preserving
embedding and lossless; compactness survives as a discrete spectrum and as a
physical observable (revival time, CHSH$\sim\xi$). The lossy operations are the
\emph{reverse direction} $\mathbb{R}\to S^1$, the spatial projections, and -- for
a readout -- the finite, non-periodic measurement window.

\section{States: the Bijective Bridge}

\textbf{[B]} In standard QM a qubit is a vector in $\mathbb{C}^2$:
\[
  |\psi\rangle = \alpha|0\rangle + \beta|1\rangle,
  \qquad |\alpha|^2 + |\beta|^2 = 1,\qquad \alpha,\beta\in\mathbb{C}.
\]
In FFGFT the same qubit is a point $(z,r,\theta)$ on the cylinder/torus:
\[
  z\in[-1,1],\quad r\in[0,1],\quad \theta\in[0,2\pi),\qquad z^2+r^2=1,
\]
with $z$ as projection onto the computation axis, $r$ as radial superposition
amplitude, $\theta$ as phase. Both representations are linked by a bridge:
\[
  \boxed{\;
  \alpha = \sqrt{\tfrac{1+z}{2}}\,e^{i\theta/2},
  \qquad
  \beta = \sqrt{\tfrac{1-z}{2}}\,e^{-i\theta/2}\;}
\]
and conversely $z = |\alpha|^2-|\beta|^2$, $r = 2|\alpha\beta|$,
$\theta = \arg(\alpha)-\arg(\beta)$. Consistency is immediate:
$|\alpha|^2+|\beta|^2 = 1$ and $z^2+r^2 = 1$; $|0\rangle$ corresponds to
$(1,0,*)$, $|1\rangle$ to $(-1,0,*)$.

\textbf{[S]} The conceptual difference: in QM $|\alpha|^2$ is a
\emph{probability} (Born rule), in FFGFT $r^2 = 1-z^2$ is an \emph{energy
density} of the radial configuration. Numerically both coincide at large
statistics; FFGFT is a deterministic geometric model whose statistical results
agree with the Born rule.

\section{Operators: Pauli Matrices as Rotations}

\textbf{[B]} The three Pauli matrices generate all single-qubit operations, with
algebra $\sigma_i\sigma_j = \delta_{ij}\mathbb{1} + i\epsilon_{ijk}\sigma_k$.
On the Bloch cylinder they correspond to three geometric rotations: $\sigma_z$
the phase rotation $\theta\to\theta+\pi$, $\sigma_x$ a $(z,r)$-rotation by $\pi$
(bit flip), $\sigma_y$ the orthogonal rotation with phase flip.

\begin{center}
\renewcommand{\arraystretch}{1.4}
{\small
\begin{tabular}{lcl}
\toprule
\textbf{Operator} & \textbf{Hilbert-space matrix} & \textbf{FFGFT geometry} \\
\midrule
$\sigma_z$ & $\begin{pmatrix}1&0\\0&-1\end{pmatrix}$ & Phase rotation $\theta\to\theta+\pi$ \\
$\sigma_x$ & $\begin{pmatrix}0&1\\1&0\end{pmatrix}$ & $(z,r)$-rotation by $\pi K_{\text{frak}}$ \\
$\sigma_y$ & $\begin{pmatrix}0&-i\\i&0\end{pmatrix}$ & Rotation with phase flip \\
$H$ & $\tfrac{1}{\sqrt2}\begin{pmatrix}1&1\\1&-1\end{pmatrix}$ & swap $z\leftrightarrow r$, $\theta\to\theta+\pi/2$ \\
$T$ & $\begin{pmatrix}1&0\\0&e^{i\pi/4}\end{pmatrix}$ & Phase rotation $\theta\to\theta+\pi/4$ \\
\bottomrule
\end{tabular}
}
\renewcommand{\arraystretch}{1.0}
\end{center}

\textbf{[K]} The \emph{only} FFGFT-specific deviation is the factor
$K_{\text{frak}}\approx0{.}9867$ in the $\sigma_x$ rotation. It follows from the
fractal dimension $D_f = 3-\xi$ (A040) and is a testable prediction: all
$\pi$-gates deviate systematically by $\Delta\approx0{.}013\,\%$ from the ideal
value. In the Hilbert-space formalism this would be an ad hoc correction to be
built in as hardware noise; in FFGFT it follows from the geometry.

\section{Gates and Universality}

\textbf{[B]} Every unitary single-qubit gate $U\in U(2)$ is a Bloch rotation,
\[
  U = e^{i\alpha}R_{\hat n}(\phi)
    = e^{i\alpha}\bigl(\cos(\phi/2)\,\mathbb{1} - i\sin(\phi/2)\,\hat n\cdot\vec\sigma\bigr),
\]
in FFGFT the rotation of the $(z,r,\theta)$ point around axis $\hat n$ by $\phi$.
The CNOT gate translates as geometric conditioning: \emph{if $z_A=-1$ (control
qubit on $|1\rangle$), apply bit flip on qubit B} -- topologically a linking of
two cylinders (entanglement). Since the universal set $\{H,T,\text{CNOT}\}$
(Solovay--Kitaev) is fully geometrically representable, FFGFT is \emph{universal
for quantum computation}: there is no QM statement not formulable in FFGFT.

\section{Tensor Products, Bell States, CHSH}

\textbf{[B]} For $n$ qubits the Hilbert space is
$\mathcal{H}_n = \bigotimes_{i=1}^n\mathbb{C}^2$ with dimension $2^n$. In FFGFT
$n$ qubits correspond to $n$ mode configurations $(z_i,r_i,\theta_i)$ on a
shared torus geometry; entanglement is a \emph{topological connection} -- knots
not reducible to product form. The Bell state $|\Phi^+\rangle =
\tfrac{1}{\sqrt2}(|00\rangle+|11\rangle)$ is a topologically connected cylinder
configuration. Both formalisms reproduce the exponential growth $\dim = 2^n$
with $2\cdot2^n-2$ real parameters.

\textbf{[K]} \textbf{CHSH prediction.} Standard QM gives the Tsirelson bound
$|\langle\text{CHSH}\rangle|\le 2\sqrt2$. FFGFT predicts a small, geometrically
derived deviation $\Delta\text{CHSH}\approx10^{-5}$ -- the \emph{elementary}
deviation per measurement. It lies below the current hardware noise floor and is
an open task for high-precision Bell tests (A160/A210). Distinct from this is the
cumulative value over many measurements; the two quantities are not directly
comparable.

\section{Measurement as Pole Transition}

\textbf{[S]} In standard QM a measurement in the computation basis gives value
$0$ with probability $|\alpha|^2$, with collapse to $|0\rangle$ or $|1\rangle$.
In FFGFT measurement is \emph{not} a probabilistic collapse but a geometric
interaction driving the $(z,r,\theta)$ point to the nearest stable extreme
($z=+1$ or $z=-1$) -- a deterministic motion in mode space that is
statistically indistinguishable from Born-rule behaviour because of the unknown
phase $\theta$. At large statistics:
\[
  P_{\text{FFGFT}}(z=+1) = \tfrac{1+z}{2} = |\alpha|^2 = P_{\text{QM}}(0).
\]
The statistical results are identical; the Bell-test debate over local realism
dissolves geometrically, because entanglement is represented as topological
connection, not as non-local correlation.

\section{Unitary Evolution: Schrödinger as Low-Energy Limit}

\textbf{[B]} The QM time evolution $i\hbar\,\partial_t|\psi\rangle = H|\psi\rangle$
is in FFGFT the deterministic motion of the mode configuration $\varphi_{n,l,j}$
on $T^3\times\mathbb{R}$. The FFGFT mode equation reduces in the low-energy limit
($E\ll E_\xi$) to the Schrödinger equation:
\[
  i\hbar\,\frac{\partial\varphi_{n,l,j}}{\partial t}
  \;\xrightarrow{E\ll E_\xi}\;
  -\frac{\hbar^2}{2m}\nabla^2\varphi_{n,l,j} + V\varphi_{n,l,j}.
\]
At higher energies $\xi$-corrections appear, which in the Hilbert-space formalism
appear as Hamiltonian modifications.

\section{What Is Translatable and What Remains FFGFT-Specific}

\textbf{[B]} \emph{All} Hilbert-space statements are mappable into FFGFT: states,
operators, gates (with or without $K_{\text{frak}}$), tensor products,
entanglement, Bell/GHZ states, measurement, unitary evolution, algorithms (Shor,
Grover, VQE). And conversely: the $(z,r,\theta)$ configuration corresponds
uniquely to the $(\alpha,\beta)$ vector, geometric rotations to unitary operators,
topological entanglements to entangled Hilbert-space vectors. Every computation
in one formalism is executable in the other, with identical result.

\textbf{[K]} The translatability holds for \emph{statements about quantum
systems}, not for the \emph{input values}. In Hilbert space $\xi$, the fermion
masses, $\alpha$, and $K_{\text{frak}}$ appear as external inputs; in FFGFT they
follow from the geometry ($\xi$ as stipulation A020; the masses from
$m_i = r_i\,\xi^{p_i}v$, A100/A110; $\alpha$ from A130; $K_{\text{frak}}$ from
A040). This is the structural asymmetry: Hilbert space has all the tools for QM
computations, but no explanation for the input values; FFGFT delivers both from
one geometry. In this sense QM is a subset of FFGFT -- restricted to the question
of \emph{how} systems behave, without the question of \emph{why} the constants
have these values.

\section{Operational Equivalence and Interpretative Divergence}

\textbf{[S]} Complete translatability must not be read as ontological identity.
\emph{Operationally} FFGFT and standard QM agree in every measurable result, up
to $\xi$-corrections. \emph{Interpretatively} they diverge at points where the
standard reading makes statements that in FFGFT appear as artefacts of the chosen
carrier assumptions. Three points within QM:

\begin{itemize}[leftmargin=2em,itemsep=3pt,topsep=3pt]
  \item \textbf{Born rule as irreducible randomness.} Standard reading: $|\psi|^2$
    as ontological probability, the outcome fundamentally random. FFGFT reading:
    same distribution, but from deterministic pole-transition dynamics on $T^4$;
    stochasticity arises at the measurement apparatus, not in the natural process.
    Operationally identical.
  \item \textbf{Non-locality as ontological property.} Standard reading of Bell
    correlations: action at a distance in $\mathbb{R}^3$. FFGFT reading:
    entanglement is a topological connection on $T^4$ -- geometric proximity in a
    closed geometry, not action at a distance in open space (A160). Results
    identical, ontological reading structurally different.
  \item \textbf{$\hbar$ as ontological primitive.} Standard reading: $\hbar$ is a
    fundamental constant of nature. FFGFT reading: $\hbar$ is an SI conversion
    factor between energy and time measures, following from the closure condition
    on $T^4$ (A080/A085). Numerically identical, ontologically derived.
\end{itemize}

\section{The Four Extensions: Hilbert Space Upward}

\textbf{[B]} The downward translation loses a structural layer at each reduction.
The reverse question is: which extensions must one \emph{add} to the standard
Hilbert space so that it natively carries the full FFGFT? Remarkably, each
extension corresponds to an established mathematical structure -- FFGFT invents no
new mathematics, but is a specific combination of existing structures.

\textbf{Extension 1 -- deformed algebra $SU_q(2)$.} The isotropic $SU(2)$ algebra
$[\sigma_i,\sigma_j]=2i\epsilon_{ijk}\sigma_k$ becomes anisotropic in FFGFT:
$\sigma_x^{\text{FFGFT}} = K_{\text{frak}}\,\sigma_x$, while $\sigma_y,\sigma_z$
are undisturbed. This is the $q$-deformation (Drinfeld--Jimbo 1985) with
\[
  \boxed{\;q = e^{i\pi(1-K_{\text{frak}})} = e^{i\pi\cdot100\xi} = e^{i\pi/75}
        \approx 0{.}999 + 0{.}042\,i\;}
\]
The deformation phase is small ($\pi/75\approx0{.}042$ rad) but structurally
significant: it breaks spherical symmetry to an axis-distinguished algebra --
the cylinder's memory of the torus. Cost: the Hilbert space remains $\mathbb{C}^2$,
the $2\times2$ matrices remain; only the structure constants change. Consequence:
all $\pi$-gates deviate by $0{.}013\,\%$ (testable).

\textbf{Extension 2 -- cyclic winding quantum number.} The cylinder has one compact
loop; the 2-torus has two. To represent the second, the Hilbert space carries a
quantum number $n\in\mathbb{Z}$: from $|\psi\rangle = \alpha|0\rangle+\beta|1\rangle$
one gets $|\psi;n\rangle = \alpha|0;n\rangle+\beta|1;n\rangle$, and
\[
  \boxed{\;\mathcal{H}^{T^2}_{\text{FFGFT}} = \bigoplus_{n\in\mathbb{Z}}\mathbb{C}^2_n\;}
\]
a countably infinite-dimensional space. The winding excitations have energy
$E_n\sim n^2\hbar^2/(mR^2)$; at low energies only $n=0$ is occupied and the space
effectively reduces to $\mathbb{C}^2$ -- the standard Bloch sphere.

\textbf{Extension 3 -- bundle structure instead of tensor product.} In standard
QM a spin-$\tfrac12$ particle is $\mathcal{H} = L^2(\mathbb{R}^3)\otimes\mathbb{C}^2$
with $[L_i,S_j]=0$; spin-orbit coupling $H_{\text{SO}}=\xi(r)\,\vec L\cdot\vec S$
must be introduced as an additional term with a free parameter. In FFGFT both
position \emph{and} spin arise from windings on the same $T^3$; the correct
Hilbert space is \emph{not} a tensor product but the section space of a spinor
bundle,
\[
  \boxed{\;\mathcal{H}^{T^3}_{\text{FFGFT}} = L^2(T^3, S)\;}
\]
with $[L_i,S_j]\neq0$ (corrections $\sim\xi$). Spin-orbit coupling is thereby a
\emph{geometric} property of the bundle, not a free parameter.

\textbf{Extension 4 -- temporal winding $k_t$.} In standard QM time is a
continuous linear axis with $E=\hbar\omega$. In FFGFT it is a compact winding
direction with $k_t\in\mathbb{Z}$ and
\[
  \boxed{\;E = \hbar f(k_t)\;}
\]
for a periodic function $f$; in the lowest sector $E\approx\hbar\omega$. The
central statement: the \emph{mass} of a mode is its $k_t$-winding energy,
$m_i c^2 = \hbar f(k_t^{(i)})$, which together with the Yukawa form
$m_i = r_i\,\xi^{p_i}v$ (A100) delivers all fermion masses from the winding
tuples.

\section{The Structural Point: Known Mathematics, New Combination}

\textbf{[B]} Each of the four extensions corresponds to a well-known structure:

\begin{center}
\renewcommand{\arraystretch}{1.4}
{\small
\begin{tabular}{clc}
\toprule
\textbf{Ext.} & \textbf{Mathematics} & \textbf{Established since} \\
\midrule
1 & Quantum groups $SU_q(2)$ & Drinfeld/Jimbo 1985 \\
2 & Winding excitations & String theory, Kac--Moody \\
3 & Spinor bundles & Differential geometry 1950s \\
4 & Compact extra dimension & Kaluza 1921, Klein 1926 \\
\bottomrule
\end{tabular}
}
\renewcommand{\arraystretch}{1.0}
\end{center}

None of these structures is new. What makes FFGFT specific is the
\textbf{combination}: all four together, with a single geometric justification
($T^4$ and $\xi$).

\section{The Full Hilbert Space}

\textbf{[B]} With all four extensions one obtains a complete FFGFT-equivalent
formulation. In the canonical A-series notation the base space is
\[
  \mathcal{H} = L^2(T^4)\otimes\mathbb{C}^3,
\]
where the $\mathbb{C}^3$ factor carries the three-element $Z_3$ identification
(A030); the full formulation supplements it with the spinor bundle $S$ and the
deformed $SU_q(2)$ algebra. In this language the FFGFT results are directly
derivable from the Hilbert-space structure: the masses from the $k_t$-windings
(A100), $\alpha$ from an eigenvalue relation (A130), spin-orbit coupling from the
bundle structure, the $K_{\text{frak}}$ corrections from the quantum-group
deformation (A040).

\textbf{[K]} On the same space the mass operator acts as a $Z_3$ circulant; from
its diagonalisation follow the Koide condition $Q=2/3$ and the angle $\theta=2/9$.
This computation is \emph{not} carried out here but belongs substantively in the
lepton sector: the operator and its diagonalisation are in A110, the angle
$\theta=2/9$ in A120. This document only provides the space on which they act.

\section{Quantum Graphity as Sketched Subset}

\textbf{[S]} The same subset strategy -- Bloch sphere $\subset$ FFGFT through
reduction -- can be applied to another established, published quantum gravity
theory: \emph{Quantum Graphity} (Konopka, Markopoulou, Severini, Smolin;
\emph{Phys.~Rev.~D}~77, 104029, 2008; arXiv:0801.0861) -- a background-independent
model in which space and geometry emerge from a dynamical weighted graph. It is
mathematically clearly defined and a concrete comparison object. It is a comparison
with established physics, not a load-bearing argument of FFGFT; and -- the crucial
difference from the Hilbert-space bijection above -- the reduction is \emph{not
carried out} but a hypothesis. Hence the consistent $[S]$ marking throughout.

Five reduction steps: (1) continuum to discrete (triangulation of $T^4$, not
proved); (2) topology forgotten ($K_N$ vs.\ $T^4$ winding structure); (3) windings
to spin labels (physical meaning of $(n_x,n_y,n_z,k_t)$ lost); (4) geometry
parameter to phase parameter ($\xi=4/30000$ fixed vs.\ free couplings
$g_V,g_B,\ldots$); (5) matter sector discarded (all 12 fermion masses absent).

What survives: one genuine correspondence -- emergent $U(1)$ from microstructure
(Levin--Wen, extended version only). FFGFT contains Hilbert-space QM provably as
a subset, and discrete-geometry programmes plausibly as a special case.

\section{Field Theory: Modes, Lagrangian, Schrödinger Limit}

\textbf{[B]} The complete field theory of the Hilbert space -- explicit mode
functions $\phi_{n,l,j}$, free Lagrangian $\mathcal L_0$, bridge formula
Lagrangian$\,\leftrightarrow\,$mode operator, Dirac line, and Schrödinger limit
-- is in A075. Here only the key formulas needed for the Hilbert-space connection.

\textbf{[B]} \textbf{Bridge formula.} Lagrangian and mode operator describe
the same physics (derivation in six steps: A075):
\[
  \int\tfrac12(\partial\,\delta m)^2\,d^4x
  = \tfrac12\langle\psi|\Phi^2|\psi\rangle,
  \qquad \Phi = \sum_i m_i P_i.
\]
Mass arises not as a Lagrangian parameter but as an eigenvalue of the mode operator.

\section{Symbol Glossary}

Symbols used throughout the entire A-series are explained in A010.
Specific to this document:

\begin{center}
\renewcommand{\arraystretch}{1.3}
{\small\begin{tabular}{lp{9cm}}
\toprule
Symbol & Meaning \\
\midrule
$q = e^{i\pi/75}$ & Deformation parameter of the Drinfeld--Jimbo quantum group;
  follows from $q = e^{i\pi\cdot100\xi}$ with $\xi=4/30000$ \\
$U_q(\mathfrak{su}(2))$ & $q$-deformed $\mathfrak{su}(2)$ algebra \\
$[n]_q$ & $q$-number: $[n]_q = (q^n-q^{-n})/(q-q^{-1})$ \\
$\sigma_x,\sigma_y,\sigma_z$ & Pauli matrices (spin operators) \\
$S_\text{CHSH}$ & CHSH value, theoretical maximum $2\sqrt2\approx2{.}828$ \\
\bottomrule
\end{tabular}}
\renewcommand{\arraystretch}{1.0}
\end{center}

\section{Verification Scripts}

\begin{itemize}[leftmargin=2em,itemsep=2pt,topsep=2pt]
  \item Bijective bridge $(z,r,\theta)\leftrightarrow(\alpha,\beta)$: numerical
    verification of $|\alpha|^2+|\beta|^2=1$ and $z^2+r^2=1$ over random states,
    plus operator translation Pauli$\,\leftrightarrow\,$rotation:
    \texttt{python/A\_Serie\_Skripte/a070\_tensorfaktor.py}
  \item $q$-deformation and factor structure: $q = e^{i\pi/75}$, limit $q\to1$,
    and the tensor factor $L^2(T^4)\otimes\mathbb{C}^3$:
    \texttt{python/A\_Serie\_Skripte/a070\_zirkulant\_dft.py}
\end{itemize}

\section{Open Questions}

\textbf{The $\Delta\text{CHSH}\sim10^{-5}$ deviation is calculated, not measured.}
The CHSH scaling formula is fully derived in A165; the computational framework
($\phi$-QFT, Bell damping, T0-Shor) and the hardware status are there. The
IBM-Heron-r2 measurement (May 2026) confirms Bell violation ($S=2{.}74$, 96{.}9\,\%
of the Tsirelson bound), but the device-to-device variation ($\sim1{.}4\,\%$)
exceeds the $\xi$ effect ($\sim0{.}001\,\%$) by a factor $\sim1400$ -- not
resolvable with today's NISQ hardware in principle (A165). As a prediction with
falsification criteria, $\Delta\text{CHSH}$ stands in A210.

\textbf{Top quark, Down-type quarks, and quark $r_i$ deep structure are unclear.}
The lepton exponents and Up-type quarks (up to Charm) follow the geometric $2/3$
sequence (A075/A100); the top quark ($p_t=-1/3$, no known geometric reason) and
the Down-type arithmetic step size remain open (A150).

\section{Supersedes}

This document subsumes: Dok.~230 (translatability FFGFT$\,\leftrightarrow\,$Hilbert-space
QM, bijection, operators, gates, measurement, interpretative divergence), Dok.~231
(the four Hilbert-space extensions), Dok.~232 (Quantum Graphity as sketched subset),
and the framework section of Dok.~282 (the Hilbert space explicit). Incorporated:
R41 (location of losses: the type-I decompactification is lossless; lossy are
the reverse direction, spatial projection, and measurement window). \emph{Not}
taken over: the mass operator/Koide part of Dok.~282 belongs substantively in
A110; the HLV/CP comparison part of Dok.~282 remains independent as an IPI
framework comparison (A240).

\section{Use of AI Tools}

This document was created with the assistance of AI tools.
Responsibility for content and errors rests entirely with the author.
Full wording of the rules in A010.
